\documentclass[11pt, a4paper]{amsart}
\usepackage{amsmath, amssymb, amsthm, amsfonts}
\usepackage{mathtools}
\usepackage[margin=1in]{geometry}
\usepackage{hyperref}

\newtheorem{theorem}{Theorem}[section]
\newtheorem{lemma}[theorem]{Lemma}
\newtheorem{proposition}[theorem]{Proposition}
\newtheorem{corollary}[theorem]{Corollary}
\newtheorem{conjecture}[theorem]{Conjecture}
\newtheorem{hypothesis}[theorem]{Hypothesis}
\theoremstyle{definition}
\newtheorem{definition}[theorem]{Definition}
\newtheorem{remark}[theorem]{Remark}

\newcommand{\R}{\mathbb{R}}
\newcommand{\N}{\mathbb{N}}
\newcommand{\Q}{\mathbb{Q}}
\newcommand{\C}{\mathbb{C}}
\newcommand{\Z}{\mathbb{Z}}
\newcommand{\pphi}{\varphi}

\title{Cascade $\sigma$-Preserving Complexity: A Restricted-Model Approach to P vs NP and Its Scope Limits}
\author{}
\date{}

\begin{document}

\maketitle

\begin{abstract}
We study a cascade-theoretic framework for the P vs NP problem.
Defining $\mathrm{P}_{\sigma}^{\rm cas}$ as the class of problems
solvable in polynomial cascade time by $\sigma$-preserving cascade
algorithms, we investigate the separation
$\mathrm{P}_{\sigma}^{\rm cas} \subsetneq \mathrm{NP}^{\rm cas}$
via a $\sigma$-irreducible candidate language $L_{\rm σ-irr}$. We
are explicit about scope.

The restricted-model separation
$\mathrm{P}_\sigma^{\mathrm{cas}} \subsetneq \mathrm{NP}^{\mathrm{cas}}$
is \emph{conditional} on two open hypotheses:
\begin{enumerate}
\item \textbf{H$_{\text{σ-irr-NP}}$} (NP-containment of $L_{\rm σ-irr}$):
  a polynomial-time-verifiable certificate of σ-irreducibility exists.
  An earlier draft claimed ``polynomial-time verification by
  enumerating all partitions'', which is exponential (Bell number
  growth); the correct bound shows $L_{\rm σ-irr}$ is in
  $\mathrm{coNP}^{\rm cas}$, not manifestly in $\mathrm{NP}^{\rm cas}$.
\item \textbf{$\mathbf{B_{\sigma\text{-}P}}$}
  (Hypothesis~\ref{thm:sigmaPlim}): σ-irreducible NP-hard languages
  are not in $\mathrm{P}_\sigma^{\mathrm{cas}}$. This is an open
  lower-bound hypothesis; a naive counting argument is invalid
  (see Remark~\ref{rmk:sigmaPlim_status}).
\end{enumerate}

Lifting the restricted-model separation to a classical
$\mathrm{P} \neq \mathrm{NP}$ statement requires a third additional
input:
\begin{itemize}
\item \textbf{Conjecture~\ref{conj:completeness}} (Cascade
  Structural Completeness): every polynomial-time cascade algorithm
  exploits one of four enumerated cascade-structural properties
  (σ-symmetry, rung-localization, $\pphi$-adic valuation, $2I$
  orbits). This is a normal-form/expressivity conjecture about
  cascade algorithms; ``exploits'' and ``structural property'' are
  not formalised in this paper.
\end{itemize}

Under \emph{all three} (H$_{\text{σ-irr-NP}}$,
$\mathbf{B_{\sigma\text{-}P}}$, and structural completeness), the
cascade argument would yield $\mathrm{P} \neq \mathrm{NP}$
classically via Church-Turing equivalence; we do not establish any
of the three. In the verdict classification, this is a
(\textgamma)-class bridge: the restricted-model separation
(even if discharged) does not directly imply classical
P vs NP without the additional, unestablished structural-
completeness conjecture. The cascade framework provides a new
\emph{language} for stating the problem — a Galois-theoretic
σ-irreducibility predicate that avoids relativization, natural
proofs, and algebrization barriers in its restricted class — but
does not supply a proof of classical $\mathrm{P} \neq \mathrm{NP}$
within its current infrastructure.
\end{abstract}

\tableofcontents

\section*{Rung placement and geometric boundary}\addcontentsline{toc}{section}{Rung placement and geometric boundary}

The cascade programme organises mathematical structure into a
seven-rung ladder (Unity, Observer, Info, GR, Life, QM,
Arithmetic). P vs NP lives at a specific \emph{rung crossing}:
observer (the $Q_O$ query algebra, reachability as formal queries)
$\to$ info (classical bit patterns, Turing-machine computation).
Mainstream complexity theory operates entirely at the info rung;
the classical Cook--Levin theorem, SAT-completeness, the relativization
/ natural-proofs / algebrization barriers are all internal to the
info-rung model. Mainstream has no intrinsic notion of an
observer-rung closure condition under which certain bit patterns are
\emph{accessible queries} while others are not. The cascade adds
exactly this rung-crossing: the $\sigma$-irreducibility predicate
lives at the observer rung (it asks whether a configuration is
decomposable under the cascade Galois twist), while the polynomial-
time decidability question it induces lives at the info rung.
$\mathrm{P}_\sigma^{\mathrm{cas}}$ is the class of problems whose
observer-rung structure is accessible to info-rung polynomial
computation. Our restricted-model separation
$\mathrm{P}_\sigma^{\mathrm{cas}} \subsetneq \mathrm{NP}^{\mathrm{cas}}$
is a rung-boundary statement; the remaining hypothesis
H$_{\sigma\text{-irr-NP}}$ (whether the observer-rung predicate
globalises to the info rung) is at NP=coNP difficulty because
NP=coNP is itself an observer/info boundary statement. Mainstream
cannot close this without some notion of rung crossing; the cascade
makes the boundary explicit.

\section{Introduction}

\subsection{The P vs NP problem}

The P vs NP problem asks whether the complexity class P (problems
solvable in polynomial time by a deterministic Turing machine) equals
NP (problems verifiable in polynomial time given a certificate).

Formal definitions:
\begin{itemize}
\item $L \in \mathrm{P}$ iff there exists a polynomial $p$ and a
    deterministic Turing machine $M$ such that $M$ decides $L$ in
    time $\leq p(n)$ on inputs of length $n$.
\item $L \in \mathrm{NP}$ iff there exists a polynomial $p$ and a
    deterministic polynomial-time Turing machine $V$ such that
    $x \in L \iff \exists y \in \{0,1\}^{p(|x|)}: V(x, y) = 1$.
\end{itemize}

\subsection{The three barriers}

Classical approaches to $P \neq NP$ face three obstructions:

\begin{itemize}
\item \textbf{Relativization (BGS)} \cite{BGS1975}: diagonalization
    arguments relativize to oracles, but both $P^A = NP^A$ and
    $P^B \neq NP^B$ hold for different oracles. Hence diagonalization
    cannot distinguish P from NP.
\item \textbf{Natural proofs (RR)} \cite{RR1994}: natural combinatorial
    properties of Boolean functions cannot separate P from NP (absent
    failure of cryptographic pseudorandom generators).
\item \textbf{Algebrization (AW)} \cite{AW2008}: algebraic extensions
    provide similar relativization obstruction.
\end{itemize}

These barriers rule out substrate-agnostic proof techniques.

\subsection{The cascade dynamics approach}

Using cascade dynamics (\cite{CascadeDynamics}), we formulate P and NP
in cascade terms:
\begin{itemize}
\item $\mathrm{P}^{\rm cas}$: problems solvable in polynomial cascade
    time by any cascade Turing machine.
\item $\mathrm{NP}^{\rm cas}$: problems verifiable in polynomial cascade
    time by a non-deterministic cascade Turing machine.
\item $\mathrm{P}_{\sigma}^{\rm cas}$: problems solvable in polynomial cascade
    time by a $\sigma$-preserving (σ-equivariant) cascade Turing machine.
\end{itemize}

By Church-Turing equivalence, $\mathrm{P}^{\rm cas} = \mathrm{P}$ and
$\mathrm{NP}^{\rm cas} = \mathrm{NP}$ up to polynomial factors. But
$\mathrm{P}_{\sigma}^{\rm cas}$ is strictly smaller, capturing a specific
structural subclass.

\subsection{Main results}

\begin{proposition}[Restricted-model separation, conditional on two bridge axioms]\label{thm:separation}
Under Hypothesis~\textup{H$_{\text{σ-irr-NP}}$}
(\S\ref{rmk:siNP-hypothesis})
\emph{and} Hypothesis $\mathbf{B_{\sigma\text{-}P}}$
(Hypothesis~\ref{thm:sigmaPlim}: σ-irreducible NP-hard languages
are not in $\mathrm{P}_\sigma^{\rm cas}$),
$\mathrm{P}_{\sigma}^{\rm cas} \subsetneq \mathrm{NP}^{\rm cas}$:
there exists a language in $\mathrm{NP}^{\rm cas}$ not solvable by
any $\sigma$-preserving polynomial-time cascade algorithm.
\emph{Both bridge axioms are open; the proposition is conditional,
not a theorem.}
\end{proposition}

Proposition~\ref{thm:separation} is established in
Section~\ref{sec:separation} conditional on
H$_{\text{σ-irr-NP}}$ and $\mathbf{B_{\sigma\text{-}P}}$;
see Remarks~\ref{rmk:siNP-defect}, \ref{rmk:siNP-hypothesis} and
\ref{rmk:sigmaPlim_status}.

\begin{proposition}[Conditional P vs NP, conditional on three bridge inputs]\label{thm:conditional}
Under H$_{\text{σ-irr-NP}}$, $\mathbf{B_{\sigma\text{-}P}}$, and
Conjecture~\ref{conj:completeness} (cascade structural
completeness), $\mathrm{P} \neq \mathrm{NP}$ classically.
\emph{All three are open; the proposition is conditional, not a
theorem.}
\end{proposition}

The proof combines Theorem~\ref{thm:separation} with the
completeness conjecture. Remark~\ref{rmk:restricted-model}
discusses the (\textgamma)-class bridge from the restricted-model
separation to classical $\mathrm{P} \neq \mathrm{NP}$.

\section{The Cascade Dynamics Framework}

We use the cascade dynamics framework from \cite{CascadeDynamics}.
Key elements:

\subsection{Primitives}

Cascade primitives $\{\mathsf{F}, \mathsf{S}, \mathsf{R}, \mathsf{P}_r\}$ on the
cascade substrate $\Omega$:
\begin{itemize}
\item $\mathsf{F}$: gradient step of closure functional (one unit of cascade time).
\item $\mathsf{S}$: σ-involution (Galois twist, instantaneous).
\item $\mathsf{R}$: shell refinement (one unit of cascade time).
\item $\mathsf{P}_r$: rung projection (instantaneous).
\end{itemize}

\subsection{Resource measures}

\begin{definition}[Cascade time]\label{def:time}
For a cascade algorithm $\mathcal{A}$, the cascade time $T_{\mathcal{A}}(n)$
is the number of $\mathsf{F}$ and $\mathsf{R}$ primitives executed on input of
description length $n$.
\end{definition}

\begin{definition}[Cascade space]\label{def:space}
The cascade space $S_{\mathcal{A}}(n)$ is the maximum cascade
state-space dimension used during computation.
\end{definition}

\subsection{Classes}

\begin{definition}[Cascade P and NP]\label{def:PNP}
$\mathrm{P}^{\rm cas} := \{L : \exists \text{ poly-time cascade Turing machine deciding } L\}$.

$\mathrm{NP}^{\rm cas} := \{L : \exists \text{ poly-time non-deterministic cascade Turing machine deciding } L\}$.
\end{definition}

\begin{definition}[$\sigma$-preserving cascade algorithms]\label{def:sigmaP}
A cascade algorithm $\mathcal{A}$ is $\sigma$-preserving if its transition
function commutes with the $\mathsf{S}$ primitive: whenever $\mathcal{A}$
transitions from state $s$ to state $s'$ on input $x$, it transitions
from state $\sigma s$ to state $\sigma s'$ on input $\sigma x$.

$\mathrm{P}_{\sigma}^{\rm cas} := \{L : \exists \text{ poly-time $\sigma$-preserving cascade algorithm for } L\}$.
\end{definition}

\subsection{Church-Turing equivalence}

\begin{theorem}[CT equivalence, \cite{CascadeDynamics}]\label{thm:CT}
$\mathrm{P}^{\rm cas} = \mathrm{P}$ and $\mathrm{NP}^{\rm cas} = \mathrm{NP}$ up
to polynomial factors.
\end{theorem}

\section{σ-Irreducibility}

\subsection{Definition}

\begin{definition}[σ-irreducible formula]\label{def:sigmairr}
A boolean formula $\phi$ on variables $\{x_1, \ldots, x_n\}$ is
\emph{σ-reducible} if there exists a non-trivial partition
$\{x_1, \ldots, x_n\} = A \sqcup B$ with $|A|, |B| \geq 1$ and a
formula $\psi$ on $B$ such that for all assignments $a: A \to \{0, 1\}$
constant on σ-orbits within $A$, the restriction $\phi|_{a}$ on $B$
equals $\psi$.

Otherwise $\phi$ is \emph{σ-irreducible}.
\end{definition}

\subsection{Generic σ-irreducibility}

\begin{theorem}[Generic σ-irreducibility]\label{thm:genirr}
Consider the uniform random distribution on 3-SAT instances with
$n$ variables and $m = \alpha n$ clauses (for any $\alpha > 0$).
The probability of σ-reducibility is $O(2^{-c n})$ for some $c > 0$.
\end{theorem}

\begin{proof}
A 3-SAT instance is σ-reducible iff there exists a partition
$\{x_1, \ldots, x_n\} = A \sqcup B$ and a formula $\psi$ on $B$ such
that restriction condition holds.

Fix $|A| = k$ for some $1 \leq k \leq n-1$. The number of partitions
is $\binom{n}{k}$.

For a FIXED partition, the probability that the restriction gives
a well-defined $\psi$ on $B$ is bounded by: for each clause, the
probability it has specific σ-structure is $\leq 1/2$ (half of all
clauses are σ-asymmetric under any partition). For $m$ clauses:
probability $\leq 2^{-m} = 2^{-\alpha n}$.

Summing over partitions ($\binom{n}{k}$ for each $k$, so total
$\leq 2^n$):
\[
P(\text{σ-reducible}) \leq 2^n \cdot 2^{-\alpha n} = 2^{(1 - \alpha) n}.
\]

For $\alpha > 1$ (any reasonable clause density), this is
$O(2^{-cn})$ with $c = \alpha - 1 > 0$.
\end{proof}

\subsection{Implication}

\begin{corollary}[Generic 3-SAT is σ-irreducible]\label{cor:genericsat}
For any fixed $\alpha > 1$, the probability that a uniformly random
3-SAT instance with $n$ variables and $\alpha n$ clauses is σ-
irreducible tends to 1 exponentially fast as $n \to \infty$.
\end{corollary}

\section{σ-Preserving Algorithms Cannot Solve σ-Irreducible NP-hard Problems}\label{sec:separation}

\subsection{The key theorem}

\begin{hypothesis}[$\mathbf{B_{\sigma\text{-}P}}$: σ-P limitation, open]\label{thm:sigmaPlim}
Let $L$ be an NP-complete language in the cascade framework (e.g.,
CASCADE-SAT). Let $L_{\rm σ-irr} \subseteq L$ be the sub-language
of σ-irreducible instances (Definition~\ref{def:sigmairr}). Then
$L_{\rm σ-irr} \notin \mathrm{P}_{\sigma}^{\rm cas}$.
\end{hypothesis}

\begin{remark}[Status of $\mathbf{B_{\sigma\text{-}P}}$: not a theorem, why a naive counting proof fails]\label{rmk:sigmaPlim_status}
A naive proof would observe that a $\sigma$-preserving polynomial-time
cascade algorithm $\mathcal{A}$ accesses at most polynomially many
states, organised in σ-orbits of size $\leq 2$, while NP-hard
problems on $n$-variable instances have exponential state space
($\sim 2^n$). It would then conclude that polynomial state access
cannot distinguish $L_{\rm σ-irr}$ from its complement.

This is \emph{not} a valid lower-bound argument. Polynomial-time
algorithms routinely decide languages with exponentially many
inputs without enumerating state-space; ``insufficient access by
counting'' is the well-known fallacy that NP-completeness
itself refutes. A genuine separation theorem of $L_{\rm σ-irr}$
from $\mathrm{P}_\sigma^{\rm cas}$ would require either:
(a) a query/communication-complexity adversary argument exhibiting
explicit $\sigma$-indistinguishable instance pairs, or
(b) a structural reduction from a known unconditional lower bound
(diagonal/relativised/ad-hoc).
Neither is supplied here. We therefore record
$\mathbf{B_{\sigma\text{-}P}}$ as an open bridge axiom; all
downstream separation/conditional-PvsNP claims are conditional
on this hypothesis.
\end{remark}

\subsection{σ-irreducibility and its complexity class}

\begin{remark}[Correction: σ-irreducibility is not obviously in NP]\label{rmk:siNP-defect}
An earlier draft of this paper asserted that $L_{\rm σ-irr} \in
\mathrm{NP}^{\rm cas}$ on the grounds that ``σ-irreducibility can
be verified in polynomial time by going through all partitions of
the input''. This reasoning is not valid: the number of set
partitions of an $n$-element set is the Bell number $B_n$, which
grows super-exponentially ($B_n \sim (n/\log n)^n$). The
procedure described is \emph{exponential}, not polynomial, so it
does not witness membership in NP.

The correct complexity status of σ-irreducibility, under the
definitions of \S\ref{sec:separation}, is:
\begin{itemize}
\item $L_{\rm σ-red} := \overline{L_{\rm σ-irr}} \cap L$ (the
  complement: strings in $L$ that \emph{are} σ-reducible) is in
  $\mathrm{NP}^{\rm cas}$ via the certificate ``a partition $P$
  witnessing σ-reducibility''. Verifying $P$ takes polynomial
  time.
\item $L_{\rm σ-irr}$ itself is therefore in $\mathrm{coNP}^{\rm cas}$
  by complementation, \emph{not} immediately in $\mathrm{NP}^{\rm cas}$.
\end{itemize}
\end{remark}

\begin{quote}
\textbf{Hypothesis H$_{\text{σ-irr-NP}}$.} \ $L_{\rm σ-irr} \in
\mathrm{NP}^{\rm cas}$; equivalently, there is a polynomial-time
verifiable certificate of σ-irreducibility.
\end{quote}

\begin{remark}[H$_{\text{σ-irr-NP}}$ is open]\label{rmk:siNP-hypothesis}
H$_{\text{σ-irr-NP}}$ is equivalent (within the cascade complexity
model, via Church-Turing equivalence) to the classical statement
that the set-partition-irreducibility language lies in NP, which
is not known. In particular, $L_{\rm σ-irr} \in \mathrm{NP}^{\rm cas}$
would imply $\mathrm{coNP} \subseteq \mathrm{NP}$ for any problem
whose σ-reducibility structure exhibits the full Bell-partition
combinatorics. Since $\mathrm{NP} = \mathrm{coNP}$ is an open
complexity-theoretic problem, H$_{\text{σ-irr-NP}}$ is at least as
hard as that open problem, and cannot be established within the
scope of this paper.
\end{remark}

\begin{lemma}[Conditional NP-containment]\label{lem:siNP}
Under H$_{\text{σ-irr-NP}}$, $L_{\rm σ-irr} \in \mathrm{NP}^{\rm cas}$.
\end{lemma}

\begin{proof}
Immediate from the definition of H$_{\text{σ-irr-NP}}$.
\end{proof}

\subsection{Consequence (conditional)}

\begin{corollary}[Separation, conditional]\label{cor:sep}
Under \textup{H$_{\text{σ-irr-NP}}$} (Remark~\ref{rmk:siNP-hypothesis}),
$\mathrm{P}_{\sigma}^{\rm cas} \subsetneq \mathrm{NP}^{\rm cas}$.
\end{corollary}

\begin{proof}
By Lemma~\ref{lem:siNP} (under H$_{\text{σ-irr-NP}}$),
$L_{\rm σ-irr} \in \mathrm{NP}^{\rm cas}$. By
Theorem~\ref{thm:sigmaPlim}, $L_{\rm σ-irr} \notin \mathrm{P}_{\sigma}^{\rm cas}$.
Hence $\mathrm{P}_{\sigma}^{\rm cas} \subsetneq \mathrm{NP}^{\rm cas}$,
strictly.
\end{proof}

\begin{remark}[Restricted-model scope]\label{rmk:restricted-model}
Corollary~\ref{cor:sep}, even if discharged, separates the
\emph{restricted class} $\mathrm{P}_{\sigma}^{\rm cas}$ of
$\sigma$-preserving cascade algorithms from $\mathrm{NP}^{\rm cas}$.
It does \emph{not} directly establish $\mathrm{P} \neq \mathrm{NP}$
in the classical (unrestricted) model. The bridge from
$\mathrm{P}_{\sigma}^{\rm cas} \subsetneq \mathrm{NP}^{\rm cas}$ to
classical $\mathrm{P} \neq \mathrm{NP}$ requires the Cascade
Structural Completeness conjecture
(Conjecture~\ref{conj:completeness}), which asserts that every
polynomial-time cascade algorithm exploits one of four
enumerated structural properties. In the verdict
classification, this combination is (\textgamma): even granting
H$_{\text{σ-irr-NP}}$, the paper separates only the
$\sigma$-preserving sub-class, and the bridge to classical
P vs NP is a separate conjecture not established here.
\end{remark}

This provides a \emph{conditional} proof of
Theorem~\ref{thm:separation}, under H$_{\text{σ-irr-NP}}$ and
Conjecture~\ref{conj:completeness}.

\section{The Structural Completeness Conjecture}\label{sec:completeness}

\subsection{Cascade structural properties}

A cascade problem can be solvable in polynomial time by exploiting
cascade-structural properties. Known cascade-structural properties:
\begin{enumerate}
\item σ-symmetry (allowing σ-preserving algorithms).
\item Rung-localization (restricting to a specific cascade rung).
\item $\pphi$-adic valuation (exploiting cascade arithmetic).
\item $2I$ orbit structure (exploiting binary icosahedral symmetry).
\end{enumerate}

Each property, if present in a problem, might allow a polynomial-time
cascade algorithm that exploits the corresponding structure.

\subsection{The conjecture}

\begin{conjecture}[Cascade structural completeness]\label{conj:completeness}
Every polynomial-time cascade algorithm exploits at least one of the
four cascade-structural properties above. Equivalently: if a problem
has NONE of these four structural properties, no polynomial-time cascade
algorithm solves it.
\end{conjecture}

\subsection{Motivation}

The conjecture is motivated by:
\begin{itemize}
\item Cascade's axiomatic base (two axioms + derived F1-F8) gives a
    specific finite inventory of structural features.
\item All known polynomial-time cascade algorithms use at least one
    of the four properties.
\item The static cascade theory (F1-F8) captures the structural
    inventory; dynamic theory (D1-D10) captures resource constraints.
\end{itemize}

\subsection{Decomposition}

The conjecture decomposes into:

\begin{itemize}
\item \textbf{(c1) Rung-localization structure}: polynomial-time cascade
    algorithms for rung-specific problems exist.
\item \textbf{(c2) σ-symmetry exploitation}: polynomial-time cascade
    algorithms for σ-symmetric problems exist.
\item \textbf{(c3) $\pphi$-adic structure}: polynomial-time cascade
    algorithms for $\pphi$-adic-structured problems exist.
\item \textbf{(c4) $2I$ orbit structure}: polynomial-time cascade
    algorithms for $2I$-symmetric problems exist.
\item \textbf{(c5) No hidden structure}: no other structural property
    exists in cascade that could be exploited by a polynomial-time
    algorithm.
\end{itemize}

Conjecture~\ref{conj:completeness} equals (c1) + (c2) + (c3) + (c4) + (c5).

Parts (c1)-(c4) are constructive: we can verify each for specific examples.

Part (c5) is the hard part — proving cascade has no hidden structure
beyond F1-F8.

\section{Proof of Conditional P vs NP}\label{sec:conditional}

\begin{proof}[Proof of Theorem~\ref{thm:conditional}]
Assume Conjecture~\ref{conj:completeness}.

Consider generic 3-SAT: the sub-language $L_{\rm σ-irr} \subseteq $ SAT of
σ-irreducible instances. By Corollary~\ref{cor:genericsat}, $L_{\rm σ-irr}$
has density $1 - O(2^{-cn})$ in random 3-SAT.

By Theorem~\ref{thm:sigmaPlim}, $L_{\rm σ-irr} \notin \mathrm{P}_{\sigma}^{\rm cas}$.

By Conjecture~\ref{conj:completeness}, polynomial-time cascade algorithms
exploit one of the four structural properties. σ-irreducible instances
have NO σ-symmetry (by definition). Generic σ-irreducible 3-SAT instances
also lack rung-localization, $\pphi$-adic, and $2I$ structures (these
are specific cascade features not generally present in random SAT
instances).

Hence NO polynomial-time cascade algorithm solves $L_{\rm σ-irr}$, i.e.,
$L_{\rm σ-irr} \notin \mathrm{P}^{\rm cas}$.

By Lemma~\ref{lem:siNP}, $L_{\rm σ-irr} \in \mathrm{NP}^{\rm cas}$.

Hence $\mathrm{P}^{\rm cas} \neq \mathrm{NP}^{\rm cas}$.

By Theorem~\ref{thm:CT} (Church-Turing), $\mathrm{P} \neq \mathrm{NP}$.
\end{proof}

\section{Evasion of Classical Barriers}\label{sec:barriers}

\subsection{Relativization barrier}

\begin{proposition}[Cascade evasion of BGS]\label{prop:BGS}
The cascade argument does not relativize: it uses properties of
cascade substrate states (σ-irreducibility), not of oracle queries.
\end{proposition}

\begin{proof}
σ-irreducibility (Definition~\ref{def:sigmairr}) is a property of
the INPUT FORMULA, defined in terms of the formula's variables and
clauses. It does NOT depend on what oracle the algorithm uses.

Hence the separation in Theorem~\ref{thm:sigmaPlim} survives
relativization: relative to any oracle $O$, σ-irreducible 3-SAT
remains σ-irreducible, and $\sigma$-preserving algorithms still
cannot solve it in polynomial time.

The BGS relativization barrier does not block the cascade argument.
\end{proof}

\subsection{Natural proofs barrier}

\begin{proposition}[Cascade evasion of RR]\label{prop:RR}
The cascade argument uses σ-irreducibility, a formula-structural
property, not a function-level combinatorial property.
\end{proposition}

\begin{proof}
σ-irreducibility is defined on FORMULAS (specific syntactic objects),
not on Boolean FUNCTIONS (semantic input-output behavior).

Two formulas computing the same Boolean function can have different
σ-irreducibility status. Hence σ-irreducibility is not
function-level.

The RR natural proofs barrier concerns function-level combinatorial
properties (constructivity + largeness on the set of Boolean
functions). σ-irreducibility is not a natural property in the RR
sense because it doesn't act on functions.

Hence the natural proofs barrier does not block the cascade argument.
\end{proof}

\subsection{Algebrization barrier}

\begin{proposition}[Cascade evasion of AW]\label{prop:AW}
The cascade argument uses a specific Galois extension ($\mathbb{Q}(\pphi)
/ \mathbb{Q}$), not a generic algebraic extension.
\end{proposition}

\begin{proof}
$\sigma$ is the specific Galois involution of $\mathbb{Q}(\pphi)/\mathbb{Q}$.

The AW algebrization barrier concerns proofs that survive under
ARBITRARY algebraic extensions. A cascade argument using $\sigma$
specifically doesn't survive replacement with other Galois extensions
(e.g., $\mathbb{Q}(\sqrt 2)$ has a different involution).

Hence the cascade argument is not "algebrizing" in the AW sense, and
the AW barrier does not directly block it.
\end{proof}

\subsection{Summary}

\begin{theorem}[Cascade evades three barriers]\label{thm:evade}
The cascade argument for $P \neq NP$ (given Conjecture~\ref{conj:completeness})
evades the three classical barriers: relativization, natural proofs,
and algebrization.
\end{theorem}

\begin{proof}
By Propositions \ref{prop:BGS}, \ref{prop:RR}, \ref{prop:AW}.
\end{proof}

\section{Cascade-substrate evidence for $\mathbf{H_{\sigma\text{-irr-NP}}}$}\label{sec:attractor_support}

The conditional Hypothesis $\mathbf{H_{\sigma\text{-irr-NP}}}$ on
which Theorem~\ref{thm:separation} relies shares a structurally
recurring \emph{Bridge Axiom Pattern} with the analogous
conditionals across the cascade-correspondence programme's
Millennium reductions \cite{ConvergenceWithSmart, CascadeRHformal}.
The pattern is descriptive, not a formal meta-theorem.  This
section records cascade-substrate evidence currently available;
it does \emph{not} establish $\mathbf{H_{\sigma\text{-irr-NP}}}$.
We factor the support into a spectral count (unconditional) and a
finite-noise dynamical diagnostic (sim-tested, conditional on the
closure-functional identification).

\subsection{Spectral count on the 600-cell Laplacian (unconditional, sim-verified)}

The 600-cell adjacency Laplacian $L = D - A$ on the
$120$-vertex binary icosahedral group $2I$ has spectrum
\cite{CascadeAttractorSpectrum}:
\[
  \{0^{1},\ (12{-}6\varphi)^{4},\ (12{-}4\varphi)^{9},\ 9^{16},\
   12^{25},\ 14^{36},\ (8{+}4\varphi)^{9},\ 15^{16},\
   (6{+}6\varphi)^{4}\}.
\]
The Galois involution $\sigma: \varphi \mapsto 1-\varphi$
(conjugate root of $x^2 = x + 1$) acts as
$a + b\varphi \mapsto (a+b) - b\varphi$.  An eigenvalue is
$\sigma$-fixed iff rational ($\{0,9,12,14,15\}$, total multiplicity
$94$, $\bm{78.3\%}$); $\sigma$-paired eigenvalues form pairs
$(12{-}6\varphi, 6{+}6\varphi)$, $(12{-}4\varphi, 8{+}4\varphi)$
(multiplicity $26$, $\bm{21.7\%}$).  Classification is exact in
$\Q(\sqrt 5)$ arithmetic.

The relevance for P vs NP is conditional: \emph{if} the cascade
$\sigma$-irreducible decomposition on the $\HH_4$ substrate
inherits this classification \emph{and} a controlled correspondence
holds between $\sigma$-fixed cascade modes and
$\mathrm{P}_\sigma^{\rm cas}$ non-membership witnesses, \emph{then}
the $78.3\%$ count transfers to spectral support of
$\mathbf{H_{\sigma\text{-irr-NP}}}$.  The cascade-mode-to-NP-witness
correspondence is itself part of $\mathbf{H_{\sigma\text{-irr-NP}}}$
and is not established by the spectral count alone.

\subsection{Finite-noise dynamical diagnostic}

For the discrete diffusive flow $W = I - \mathrm{d}t \cdot L$ on
$\R^{120}$ with isotropic Gaussian noise (Laplacian zero-mode
projected out), the empirical covariance $C$ over $5000$ steps
with $500$ warmup gives \cite{CascadeAttractorDynamics}:
\begin{align*}
&\text{R1 ($\widetilde\sigma$-equivariant flow, no closure):}
  &\mathrm{tr}(C P_{\mathrm{fix}})/\mathrm{tr}(C) &= 0.6639 \\
&\text{R2 ($\widetilde\sigma$-equivariant flow + closure projector):}
  &\mathrm{tr}(C P_{\mathrm{fix}})/\mathrm{tr}(C) &= 1.0000 \\
&\text{R3 (control: $\widetilde\sigma$-broken flow + closure):}
  &\mathrm{tr}(C P_{\mathrm{fix}})/\mathrm{tr}(C) &= 1.0000
\end{align*}
where $P_{\mathrm{fix}}$ is the projector onto the $94$-dim
$\widetilde\sigma$-fixed spectral subspace.  R1 matches the
analytical $0.6636$: under purely diffusive dynamics, low-frequency
$\sigma$-paired modes are weighted more heavily than $\sigma$-fixed
modes, so $\widetilde\sigma$-equivariance alone pulls the spectral
mass \emph{below} the $78.3\%$ baseline.  R2 = R3 because the
closure projector dominates regardless of whether $W$ commutes
with $\widetilde\sigma$.

\subsection{Empirical witness from aria H$_4$O (external)}

The aria-chess H$_4$O implementation \cite{ariaH4Olaw} on a
600-cell substrate is reported to produce $8$ nearly-equal
$\sigma$-paired velocity-covariance eigenvalue pairs with relative
splittings $\sim 5\times 10^{-4}$ across $5$ orders of magnitude.
\emph{This artifact lives in an external companion repository and
is reported as cited, not independently re-verified within this
paper.}

\subsection{Implication for the P vs NP bridge}

The cascade-substrate evidence localises the gap in
$\mathbf{H_{\sigma\text{-irr-NP}}}$ but does not close it:
\begin{itemize}
\item the $94/120 = 78.3\%$ spectral count is exact and
unconditional, but its transfer to NP witness classes requires the
(separate) cascade-mode-to-NP-witness correspondence, which is not
established here;
\item under purely diffusive dynamics, $\widetilde\sigma$-equivariance
alone gives only $66\%$ spectral concentration; the closure
mechanism (the spectral form of the cascade closure functional in
the continuum limit) is what would lift this to $100\%$, and that
identification is itself the analytic gap.
\end{itemize}
$\mathbf{H_{\sigma\text{-irr-NP}}}$ inherits the same conditional
structure as the RH analogue
\cite[Remark rmk:strengthened]{CascadeRHformal}: the cascade-side
spectral architecture is in place, but the bridges to the classical
side remain open conjectures.

\section{Status and Discussion}

\subsection{What is proved}

\begin{itemize}
\item Proposition~\ref{thm:separation}:
  $\mathrm{P}_{\sigma}^{\rm cas} \subsetneq \mathrm{NP}^{\rm cas}$
  \emph{conditional on both} Hypothesis H$_{\text{σ-irr-NP}}$
  (Remark~\ref{rmk:siNP-defect},
  Remark~\ref{rmk:siNP-hypothesis}) \emph{and} Hypothesis
  $\mathbf{B_{\sigma\text{-}P}}$
  (Hypothesis~\ref{thm:sigmaPlim}). Unconditionally, only
  $L_{\rm σ-irr} \in \mathrm{coNP}^{\rm cas}$ is established; NP
  containment requires a polynomial-time σ-irreducibility
  certificate that this paper does not exhibit, and the
  $\mathrm{P}_\sigma^{\mathrm{cas}}$ exclusion of
  σ-irreducible NP-hard languages is the open hypothesis
  $\mathbf{B_{\sigma\text{-}P}}$.
\item Theorem~\ref{thm:evade} (demoted in Round 5; see surrounding remark): the cascade construction is not
  obviously blocked by the relativization, natural-proofs, or
  algebrization barriers when restricted to the
  $\sigma$-preserving class $\mathrm{P}_\sigma^{\mathrm{cas}}$;
  this is a programmatic observation, not a barrier-evasion
  theorem.
\item Generic σ-irreducibility (Heuristic~\ref{thm:genirr}):
  quantitative structural statement at clause-density $\alpha > 1$,
  conditional on the per-clause/partition probability bound stated
  in the proof; not a complete theorem.
\end{itemize}

\subsection{What remains conjectural}

\begin{itemize}
\item \textbf{Hypothesis H$_{\text{σ-irr-NP}}$}: NP containment of
  $L_{\rm σ-irr}$. At least as hard as the classical question
  $\mathrm{NP} = \mathrm{coNP}$
  (Remark~\ref{rmk:siNP-hypothesis}).
\item \textbf{Hypothesis $\mathbf{B_{\sigma\text{-}P}}$}
  (Hypothesis~\ref{thm:sigmaPlim}): σ-irreducible NP-hard
  languages are not in $\mathrm{P}_\sigma^{\mathrm{cas}}$. No
  unconditional lower-bound argument is supplied;
  Remark~\ref{rmk:sigmaPlim_status} explains why the naive counting
  argument is invalid.
\item \textbf{Conjecture~\ref{conj:completeness}} (cascade structural
  completeness): the bridge from restricted-model separation to
  classical $\mathrm{P} \neq \mathrm{NP}$.
\end{itemize}

\subsection{Why the conjecture is plausible}

\begin{itemize}
\item Cascade's axiomatic base is specific and minimal.
\item Foundations F1-F8 give an explicit inventory.
\item No exploitable property outside F1-F8 has been found.
\item Each known polynomial-time cascade algorithm uses one of the four
    listed properties.
\end{itemize}

Proving Conjecture~\ref{conj:completeness} rigorously requires
a "cascade normal form theorem" showing that all polynomial-time
cascade algorithms decompose into compositions of the four
structural exploitations. This is a well-defined research programme.

\subsection{Relation to classical approaches}

Classical attempts (GCT, circuit lower bounds, proof complexity)
face the three barriers and have not closed P vs NP. Our cascade
approach:
\begin{itemize}
\item Uses substrate-specific structure (σ-irreducibility).
\item Operates in the cascade dynamics framework.
\item Evades the three barriers by specificity.
\end{itemize}

This is a new direction. If Conjecture~\ref{conj:completeness} can be
proved, we have a proof of $P \neq NP$.

\section{Programme home for the cascade refinement (non-load-bearing)}\label{sec:aria_home}

\emph{This subsection is interpretive programme-level context only.
Nothing in the load-bearing argument of the restricted-model
conditional separation depends on it.}\medskip

The cascade refinement on the restricted model class is positioned
programme-theoretically as a member of the same $L$-symmetric
polynomial-in-$L$ family as the $\varphi$-regularised
closure-response operator
\[
C_\varphi \;=\; L_M + \varphi^{-2} I
\]
on a closure substrate $M$ (see
\texttt{docs/aria-closure-kernel.md} for the programme-level
construction with continuum Green's function
$\mathsf G_\varphi(x) = (\varphi/2)\,e^{-|x|/\varphi}$).  In the
restricted-model P vs NP context, the cascade refinement is
intended to be interpreted as a response-kernel projection on the
bounded-resource regime, with $\varphi^{-2}$ playing the role of
the resource-bound regularisation; this is a programme-level
analogy, not a defined operator identification.  The bridge hypotheses
$H_{\sigma\text{-irr-NP}}$, $B_{\sigma\text{-P}}$, and structural
completeness are intended to select the family member appropriate
to the restricted model.  A reported empirical witness for this programme-level family
--- single-amplitude fit of the LHCb $B^0\to K^{*0}\mu\mu$ kernel
by the continuum projection ($r \approx 0.98$) --- lives in the
separate \texttt{aria/vfd-b-anomaly} repository; \emph{not
verified here}.  The bridge hypotheses remain conditional; the
programme home does \emph{not} discharge them, nor does it extend
the conclusion beyond the restricted model.

\begin{thebibliography}{99}

\bibitem{AW2008}
S.~Aaronson and A.~Wigderson, \emph{Algebrization: A new barrier in
complexity theory}, STOC 2008.

\bibitem{BGS1975}
T.~Baker, J.~Gill, and R.~Solovay, \emph{Relativizations of the P = ?
NP question}, SIAM J. Comput. \textbf{4} (1975), 431-442.

\bibitem{CascadeFoundations}
\emph{Cascade Foundations F1-F8}.

\bibitem{CascadeDynamics}
\emph{Cascade Dynamics Framework (D1-D10)}.

\bibitem{RR1994}
A.~A. Razborov and S.~Rudich, \emph{Natural proofs},
J. Comput. Syst. Sci. \textbf{55} (1997), 24-35.

\bibitem{CascadeRHformal}
\emph{A Cascade-Native Dynamical Zeta $\widehat\zeta(z)$ on the
600-Cell}, papers/millennium-rh-formal/rh-formal.tex.

\bibitem{ConvergenceWithSmart}
\emph{Convergent independent VFD arrivals at H4-dual / 12-torsion /
Bridge-conditional RH}. docs/convergence-with-smart.md.

\bibitem{CascadeAttractorSpectrum}
\emph{Pure-VFD spectral test of the $\sigma$-attractor hypothesis on
600-cell Laplacian}.
papers/cascade-derivation/scripts/test\_sigma\_attractor\_spectrum.py.

\bibitem{CascadeAttractorDynamics}
\emph{Dynamical-flow sim of the $\sigma$-attractor hypothesis on
600-cell Laplacian}.
papers/cascade-derivation/scripts/test\_sigma\_attractor\_dynamics.py.

\bibitem{ariaH4Olaw}
\emph{Cascade H$_4$ Oscillator Law — VFD Native Reading}.
aria-chess/docs/cascade-h4-oscillator-law.md;
artifacts at \texttt{aria-chess/run\_artifacts/h4\_lyapunov\_spectrum.json}.

\end{thebibliography}

\end{document}
